% sec-01-fda-forms.tex - FDA Forms 1571 and 3674 (full stage)
\section{FDA Forms 1571 and 3674}

The administrative gateway to this initial Investigational New Drug
(IND) application is the pair of executed Food and Drug Administration
(FDA) forms that open every Center for Drug Evaluation and Research (CDER)
submission: FDA Form 1571 (Investigational New Drug Application) and FDA
Form 3674 (Certification of Compliance, under 42 U.S.C. \S 282(j)). These
two forms are positioned, together with the Cover Letter, ahead of the
Table of Contents so that a reviewer encounters the sponsor agreement, the
serial number, the contents map, and the clinicaltrials.gov certification
before any scientific module. They are the regulatory envelope around the
LLM-Directed PDAC Robotic Daraxonrasib Phase 1 protocol (protocol
v1.0.0), a first-in-human, prospective, open-label, single-arm,
early-feasibility study that combines an IND under 21 CFR Part 312 with an
Investigational Device Exemption (IDE) framework under 21 CFR Part 812 for
an on-premises, LLM-directed, eight-arm robotic pancreaticoduodenectomy
(Whipple) device used together with oral daraxonrasib (RMC-6236), an oral
RAS(ON) multi-selective (pan-RAS) inhibitor \cite{daraxwhipple,
onpremwhippl}. This section summarizes the field-by-field content and the
governing logic of each form; the executed, signed PDF versions of FDA
Form 1571 and FDA Form 3674 are assembled with the physical submission and
are not reproduced inside this LaTeX source, in keeping with FDA's
expectation that the original signed forms accompany the package
\cite{phase1guidance, cfr312paiuni}.

Figure 1 traces the cover-sheet flow from the Cover Letter through Form
1571 (its serial logic and dual purpose) into Form 3674 (its three
mutually exclusive certification boxes) and onward to the assembled IND
that starts the 30-day FDA review clock under 21 CFR \S 312.40.

\begin{mermaidfig}
\node[mmaccent] (cl)   at (0,0)      {Cover Letter\\sponsor-investigator,\\drug, serial 0000};
\node[mmgoal]   (f1)   at (0,-2.4)   {FDA Form 1571\\(1) sponsor agreement\\(2) cover sheet for\\every submission};
\node[mmin]     (ser)  at (-4.6,-4.9) {Box 10 serial number\\initial = 0000\\(0001, 0002 ... follow)};
\node[mmproc]   (f3)   at (3.6,-4.9)  {FDA Form 3674\\clinicaltrials.gov\\certification of compliance};
\node[mmctx]    (ba)   at (0.0,-7.6)  {Box A\\no clinical\\investigation};
\node[mmdec]    (bb)   at (4.2,-7.9)  {Box B\\402(j) PHS Act\\does NOT apply\\(typical Phase 1)};
\node[mmctx]    (bc)   at (8.2,-7.6)  {Box C\\402(j) applies;\\will comply};
\node[mmgoal]   (sub)  at (-1.6,-10.6) {Assembled IND\\submitted to FDA;\\30-day review\\clock starts};
\draw[mmedge] (cl) -- (f1);
\draw[mmedge] (f1) to[out=200,in=70,looseness=0.7] (ser.north);
\draw[mmedge] (f1) to[out=340,in=110,looseness=0.7] (f3.north);
\draw[mmedge] (f3) to[out=250,in=70,looseness=0.6] (ba.north);
\draw[mmedge] (f3) -- (bb.north);
\draw[mmedge] (f3) to[out=300,in=110,looseness=0.6] (bc.north);
\draw[mmedge] (bb) to[out=250,in=20,looseness=0.7] node[mmlabel]{selected for this trial} (sub.east);
\draw[mmedge] (ser) to[out=270,in=160,looseness=0.7] (sub.west);
\end{mermaidfig}
\figcaption{Figure 1. The Cover Letter, FDA Form 1571, and FDA Form 3674
cover-sheet flow. Form 1571 serves the dual purpose of a
sponsor-investigator agreement and a cover sheet for every submission; the
initial serial number in box 10 is 0000. Form 3674 certifies
clinicaltrials.gov compliance through one of Box A (no clinical
investigation), Box B (an investigation to which 42 U.S.C. \S 282(j) /
section 402(j) of the PHS Act does not apply, the typical Phase 1
selection), or Box C (the requirements apply). The Cover Letter and the
FDA forms precede the Table of Contents.}

\subsection{FDA Form 1571 - Investigational New Drug Application}

FDA Form 1571 serves a dual purpose that defines its central place in the
submission. First, it is the binding sponsor-investigator agreement: by
signing box 13, the sponsor-investigator, ChemicalQDevice (Kevin Kawchak,
Chief Executive Officer, as the named sponsor-investigator), commits that
the investigation will not begin until 30 days after FDA receives the IND
(unless notified sooner that the studies may proceed), that a reviewing
Institutional Review Board (IRB) compliant with 21 CFR Part 56 will be
responsible for the initial and continuing review of each study, that the
investigation will be conducted in accordance with all applicable
regulatory requirements, and that an investigator will obtain informed
consent under 21 CFR Part 50 \cite{cfr050paiuni, ichpaistduni}. Second,
and on a recurring basis, Form 1571 is the cover sheet that accompanies
every subsequent submission to the IND file, including protocol
amendments, information amendments, IND safety reports, annual reports,
and general correspondence; each such submission carries its own Form 1571
bearing the next sequential serial number so that the agency can
reconstruct the complete history of the file \cite{phase1guidance}.

The form opens with sponsor identification and contact fields. Box 1
(Name of Sponsor) and boxes 2 through 5 (Address, including street, city,
state, and ZIP) carry the sponsor-investigator entity ChemicalQDevice and
its San Diego, California place of business; box 6 (Telephone Number)
and the associated electronic-contact information complete the routing so
that the assigned project manager and the safety contact can be reached.
Box 7 captures the name(s) of any contract research organization, which
for this single-academic-site study is limited to the named clinical site
and the entities providing investigational product and device support.
The IND Number field is left blank on the initial submission because the
number is assigned by FDA upon receipt; once assigned, that pre-assigned
or receipt-assigned IND number is carried on the Form 1571 of every later
serial submission. Box 9 (Name of Drug) records daraxonrasib (RMC-6236),
the oral RAS(ON) multi-selective small-molecule inhibitor that binds the
GTP-bound active state of mutant RAS in complex with cyclophilin A and
suppresses downstream effector signaling, administered once daily
\cite{daraxwhipple, pdac060s2030}.

The general-information block describes the proposed research and the
phase of investigation. The form indicates that this is a Phase 1
investigation (the first-in-human checkbox), names the indication as
histologically or cytologically confirmed resectable or
borderline-resectable pancreatic ductal adenocarcinoma (PDAC) with a
documented KRAS G12 mutation (G12D, G12V, or G12R) in patients who are
candidates for pancreaticoduodenectomy, and lists the protocol title and
the combined IND and IDE nature of the application. The contents-of-the-
application checkboxes (box 12) are marked to indicate which items of 21
CFR \S 312.23 are included in this serial submission: the Table of
Contents, the Introductory Statement and General Investigational Plan, the
Investigator's Brochure, the study protocol(s), Chemistry, Manufacturing,
and Controls (CMC) information, pharmacology and toxicology information,
previous human experience, and additional information. Box 11 designates
the individuals responsible for monitoring the conduct and progress of the
clinical investigation and the person(s) responsible for review and
evaluation of information relevant to the safety of the drug; for this
study those responsibilities map onto the three-tier oversight structure
(the Data and Safety Monitoring Board, the Independent Safety Monitor, and
the Physical AI Safety Review Committee) described in the protocol
\cite{paiautospons, clintrustpai}.

Box 10 is the serial number. FDA requires every submission to an IND
to be numbered consecutively, and the initial submission is serial number
0000. Each subsequent submission to this IND file is numbered
sequentially as 0001, 0002, 0003, and so forth, in the chronological order
in which the agency receives it, regardless of the submission type. For
this initial IND the serial number is therefore 0000, and that value
appears both in box 10 of Form 1571 and in the Cover Letter so that a
reviewer can immediately confirm that this is the originating submission
that starts the 30-day review period under 21 CFR \S 312.40
\cite{phase1guidance, cfr312paiuni}. Table 1 summarizes the principal
Form 1571 fields and the value or status of each for this initial
submission.

\begin{center}
\begin{xltabular}{\textwidth}{>{\raggedright\arraybackslash}p{2.6cm} Y >{\raggedright\arraybackslash}p{3.4cm}}
\toprule
\textbf{Form 1571 field} & \textbf{Description and content} & \textbf{Value / status (serial 0000)} \\
\midrule
\endfirsthead
\toprule
\textbf{Form 1571 field} & \textbf{Description and content} & \textbf{Value / status (serial 0000)} \\
\midrule
\endhead
\bottomrule
\endfoot
Boxes 1 to 6: Sponsor and contact & Sponsor-investigator name, mailing address, and telephone for the responsible entity. & ChemicalQDevice; San Diego, California; sponsor-investigator Kevin Kawchak, CEO. \\
IND Number & Assigned by FDA on receipt; carried on every later serial submission. & Blank on initial submission (to be assigned). \\
Box 7: CRO / site & Contract research organizations and clinical site(s) supporting the study. & Single academic clinical site; no separate CRO. \\
Box 9: Name of Drug & Investigational drug and active ingredient. & Daraxonrasib (RMC-6236), oral RAS(ON) pan-RAS inhibitor. \\
Phase / indication & Phase of investigation and the disease studied. & Phase 1, first-in-human; resectable or borderline-resectable KRAS G12 PDAC. \\
Box 10: Serial Number & Consecutive numbering of every submission to the file. & 0000 (initial); 0001, 0002 ... follow. \\
Box 11: Monitoring / safety & Persons responsible for monitoring conduct and for safety review. & DSMB, Independent Safety Monitor, Physical AI Safety Review Committee. \\
Box 12: Contents & Checkboxes for the 21 CFR \S 312.23 items included. & Introductory statement, general plan, IB, protocol, CMC, pharm/tox, prior human experience, additional information. \\
Box 13: Sponsor agreement & Signed commitments (30-day wait, IRB review, GCP, informed consent). & Signed by sponsor-investigator; executed PDF assembled with submission. \\
\bottomrule
\end{xltabular}
\end{center}
\figcaption{Table 1. Field-by-field summary of FDA Form 1571 for the
initial IND (serial number 0000). The executed and signed PDF of the form
is assembled with the physical submission and is not reproduced in this
source.}

The dual purpose of Form 1571 has a practical consequence for the
lifecycle of this Physical AI study. Because each protocol amendment,
information amendment, and IND safety report under 21 CFR \S 312.32 will
itself be filed under a new serial number on its own Form 1571, the form
becomes the running index of the entire investigational record, including
the Physical AI specific reporting events. The protocol defines six
Physical AI reporting triggers under 21 CFR \S 312.32(g) (a serious
Physical AI adverse event reported on the 7-day or 15-day clock; a
system-drug interaction within 15 days; a cybersecurity incident within 15
days, or within 7 days where there is potential for patient harm;
sustained AI model degradation over at least three consecutive procedures
or 24 cumulative hours; sim-to-real divergence exceeding twice tolerance,
that is a trajectory deviation greater than 4 mm or a force deviation
greater than 1.0 N; and a digital-twin discrepancy), and each such report
will enter the file as a serially numbered submission cover-sheeted by
Form 1571 \cite{paiautospons, congress9510}. In this way the form binds
the administrative numbering scheme to the safety-surveillance scheme of
the trial.

\subsection{FDA Form 3674 - Certification of Compliance}

FDA Form 3674 is the certification, required by section 402(j) of the
Public Health Service (PHS) Act (codified at 42 U.S.C. \S 282(j)), that
the sponsor has complied with the registration and results-reporting
requirements of clinicaltrials.gov to the extent those requirements apply
to the submission. The form presents three mutually exclusive
certification boxes, exactly one of which must be selected, and the
selection determines what the sponsor is affirming about the
clinicaltrials.gov status of the studies referenced in the application
\cite{phase1guidance}. The logic of the three boxes is summarized in
Table 2.

\begin{center}
\begin{xltabular}{\textwidth}{>{\raggedright\arraybackslash}p{1.6cm} Y >{\raggedright\arraybackslash}p{3.4cm}}
\toprule
\textbf{Box} & \textbf{What the box certifies} & \textbf{Applicability to this Phase 1} \\
\midrule
\endfirsthead
\toprule
\textbf{Box} & \textbf{What the box certifies} & \textbf{Applicability to this Phase 1} \\
\midrule
\endhead
\bottomrule
\endfoot
Box A & The submission does not reference any clinical trial; that is, there is no clinical investigation at all. & Not selected. This submission references a Phase 1 clinical investigation. \\
Box B & The submission references a clinical trial, but the requirements of 42 U.S.C. \S 282(j) (section 402(j) of the PHS Act) do not apply to that trial. & Selected. This is a Phase 1 investigation; a Phase 1 trial of a drug is excluded from the "applicable clinical trial" definition that triggers mandatory 402(j) registration and results reporting. \\
Box C & The submission references a clinical trial that is an applicable clinical trial to which 42 U.S.C. \S 282(j) does apply, and the sponsor certifies it has complied (or will comply) with registration and results reporting. & Not selected for this initial Phase 1. Would apply to a later applicable (for example, Phase 2 or later controlled) trial. \\
\bottomrule
\end{xltabular}
\end{center}
\figcaption{Table 2. The Box A / B / C certification logic of FDA Form
3674 and the selection for this initial Phase 1 IND. Exactly one box is
checked; Box B is selected with the rationale stated in the text.}

For this initial Phase 1 IND, Box B is selected. The rationale is that
the study referenced in this application is a Phase 1, first-in-human
investigation of daraxonrasib together with the LLM-directed robotic
Whipple device, and a Phase 1 trial of a drug falls outside the
statutory definition of an "applicable clinical trial" that triggers the
mandatory clinicaltrials.gov registration and results-reporting
obligations of 42 U.S.C. \S 282(j). Under that definition, the
controlled clinical investigations of a drug or biologic that are subject
to mandatory registration are generally those other than Phase 1 trials;
because the present study is Phase 1 (up to n = 18 treated subjects, with
approximately 36 screened, at a single academic site, using a 3+3 dose
escalation across dose levels of 160 mg, 220 mg, and 300 mg once daily
with a 28-day dose-limiting-toxicity window), the section 402(j)
requirements do not apply to it, and Box B is therefore the correct and
typical Phase 1 selection \cite{phase1guidance, oncotrialllm}.

Selecting Box B does not signify that the trial will go unregistered.
The sponsor records that, notwithstanding the inapplicability of the
mandatory 402(j) requirements, voluntary registration on
clinicaltrials.gov is intended and is encouraged as a matter of
transparency, and that any human clinical trial supported in whole or in
part by National Institutes of Health (NIH) funding is subject to the NIH
policy on clinical-trial registration and results information submission,
which independently requires registration on clinicaltrials.gov
regardless of the Form 3674 box selected. The first-in-human Physical AI
nature of this study, and the public-trust considerations attached to an
LLM-directed surgical device operating under verification-before-
generation safeguards aligned with H.R. 9510 (the Verification Before
Generation in Physical AI Oncology Trials Act of 2026), further support
voluntary, prospective registration so that the design, the staged
autonomy limits, and the safety endpoints are publicly visible before the
first dose \cite{congress9510, clintrustpai, paipredict4x}. Should the
program advance to a later applicable clinical trial (for example, a
Phase 2 study), Box C would then be selected on that submission's Form
3674 and the corresponding clinicaltrials.gov registration and results
reporting would be certified.

As with Form 1571, the certification on Form 3674 is provided on the
executed, signed PDF that is assembled with the physical submission; the
signature attests that the information on the certification is true and
complete and acknowledges that a knowing and willful false certification
is a criminal offense under 18 U.S.C. \S 1001. Because Form 3674 must
accompany an IND submission that references one or more clinical trials,
a current Form 3674 is filed with the initial IND (serial 0000) and is
re-filed, with the box selection re-evaluated, whenever a later
submission adds or materially changes a referenced clinical
investigation, keeping the clinicaltrials.gov certification synchronized
with the evolving scope of the program \cite{phase1guidance,
paisitedocpk}. Together, the executed Form 1571 and Form 3674, placed
ahead of the Table of Contents with the Cover Letter, complete the
administrative envelope that initiates FDA's review of this LLM-directed
PDAC robotic daraxonrasib Phase 1 application.
